\documentclass[11pt,letterpaper]{article}
\usepackage[margin=1in]{geometry}
\usepackage[T1]{fontenc}
\usepackage{lmodern,microtype,amsmath,amssymb,amsthm,mathtools}
\usepackage[hidelinks]{hyperref}
\setlength{\emergencystretch}{2em}
\newtheorem{theorem}{Theorem}
\newtheorem{lemma}[theorem]{Lemma}
\newtheorem{proposition}[theorem]{Proposition}
\DeclareMathOperator{\Ric}{Ric}
\DeclareMathOperator{\Rm}{Rm}
\newcommand{\Yrd}{Y_{\mathrm{rd}}}
\newcommand{\F}{\mathcal F}
\newcommand{\dV}{\,dV}
\title{Attainment of subcritical chiral Weyl energy\\on a fixed smooth homotopy four-sphere}
\author{Edward B. Baker III}
\date{Working draft, 26 November 2026}
\begin{document}
\maketitle
\begin{abstract}
On a smooth homotopy four-sphere admitting a positive Yamabe class of chiral
Weyl energy below \(8\pi^2\), we obtain a smooth positive Bach-flat class
realizing the fixed-type infimum. The argument uses Bour's fixed-parameter
dichotomy and the marked spherical-cut results. Starting within \(\alpha^2\)
of the infimum gives an estimate for the physical trace-free Ricci integral
throughout the original flow. This excludes every multi-ended singularity
model. In the one-ended case a comparison on the entire punctured core
identifies its smooth compactification with the source. The remaining
parameter limit uses a uniformly elliptic equation for Ricci at
constant-scalar critical metrics. Positive subcritical entry and rigidity
of a minimizing class are separate questions.
\end{abstract}

\section{The variational problem}
Let \(M\) be a smooth closed connected oriented homotopy four-sphere,
with smooth Riemannian metrics. Put
\[
 E_\pm(g)=\int_M |W^\pm_g|_{\rm HS}^2\dV_g,\qquad
 I=I^+(M)=\inf_{Y([g])>0}E_+(g),\qquad T=4\pi^2.
\]
The norms are the Hilbert--Schmidt norms of the curvature endomorphisms
on \(\Lambda^\pm\). The Yamabe quotient uses \(6|du|^2+Ru^2\) over
\(\|u\|_4^2\) after integration in the numerator. The unit sphere
has positive sectional curvature and \(\Yrd=8\sqrt6\,\pi\).
On \(M\) the signature formula gives \(E_+=E_-=e\). On domains we
keep the two integrals separately. All diffeomorphisms identifying
oriented manifolds in this argument preserve the chosen orientations.

\begin{theorem}\label{thm:main}
If \(M\) admits a smooth class with \(Y>0\) and \(e<2T\), there is a
smooth class on \(M\) with \(Y>0\), Bach tensor zero and \(e=I\).
On a nonstandard smooth type \(I=a_M\), the least positive subcritical
Bach-flat value. Its Weyl second variation is nonnegative along
smooth metric variations at any metric in a minimizing class.
\end{theorem}

The standard smooth type has the round metric. Our interest is therefore
in a putative nonstandard type. The hypothesis gives \(I<2T\).
We can choose any sequence \(0<\alpha_j\downarrow0\) and unit-volume
Yamabe representatives in positive classes with
\(0\le\epsilon_j=e_{j,0}-I\le\alpha_j^2\). The inequality
\(Y([g])\le\Yrd\) and existence of smooth Yamabe representatives have
their usual compact four-dimensional hypotheses \cite{LeeParker1987Yamabe}.

\section{An estimate along the original trajectory}
Write \(D=\int|\Ric_0|^2\dV\), \(S=\int R^2\dV\), using the ordinary
two-tensor norm. Gauss--Bonnet and Bour's functional give
\[
 16\pi^2=2e+S/24-D/2,\qquad
 \F^\alpha=2(1-\alpha)e+\alpha D/2
           =2e+\alpha S/24-16\alpha\pi^2 .
\]
Let \(g(t)\) solve \(\partial_tg=-2\operatorname{grad}\F^\alpha\)
from one chosen initial representative. Set \(Y_0=Y([g(0)])\).
On this representative \(S_0=Y_0^2\) and
\(D_0=4e_0-(\Yrd^2-Y_0^2)/12\le4e_0\).

\begin{proposition}\label{prop:budget}
There are \(y>0\), \(\eta>0\) and \(j_0\) depending on the
entered sequence such that, for \(j\ge j_0\) and every smooth time on
the original maximal trajectory,
\[
 \begin{gathered}
 Y([g(t)])\ge y,\qquad I\le e(t)\le e_0/(1-\alpha)\le 2T-\eta,\\
 S(t)\le Y_0^2+48\epsilon/\alpha,\qquad
 D(t)\le 4I+4\epsilon/\alpha\le4I+4\alpha<32\pi^2 .
 \end{gathered}
\]
The last strict bound has a uniform positive margin after increasing
\(j_0\).
\end{proposition}
\begin{proof}
Put \(m=(1-\alpha)16\pi^2-\F^\alpha(g(0))
=16\pi^2-2e_0-\alpha Y_0^2/24\). Strict pinching is equivalent to
\(\alpha Y_0^2<48(2T-e_0)\). In particular \(m\) tends to
\(2(2T-I)>0\). The functional decreases. Since \(\F^\alpha\ge2(1-\alpha)e\),
the physical energy is at most \(2T-m/[2(1-\alpha)]\). In a
unit-volume Yamabe representative of the current class, Gauss--Bonnet
gives \(Y^2=48(2T-e)+12D_{\rm Yam}\), hence \(Y^2\ge24m/(1-\alpha)\).
The sign of \(Y\) agrees with the sign of the first eigenvalue of
the conformal Laplacian, which is continuous along a smooth path of
compact metrics. The initial sign is positive, so \(Y\ge\sqrt{24m}\)
throughout. This also verifies the strict application of
\cite[Lemma 4.5]{Bour2010FourthOrderCurvatureFlows}.
Thus every time slice is admissible on the original \(M\), and \(e(t)\ge I\).
Since \(\F^\alpha(g(0))\le2e_0\), the first form of \(\F^\alpha\)
gives the upper chiral bound. The second gives
\(S(t)-S_0\le48(e_0-e(t))/\alpha\).
Finally the first form gives
\[
 D(t)\le D_0+\frac{4(1-\alpha)}{\alpha}(e_0-e(t))
 \le 4e_0+\frac{4(1-\alpha)}{\alpha}\epsilon
 =4I+\frac{4\epsilon}{\alpha}.
\]
The full curvature integral in curvature-operator norm is
\(16\pi^2+D(t)\), also uniformly bounded.
\end{proof}

\section{Fixed positive regularization}
We use the strictly pinched case of
\cite[Theorems 1.4, 1.9, 1.11 and \S8]{Bour2010FourthOrderCurvatureFlows}.
We record here the proof inputs that determine its scope.
Short-time existence and uniqueness on compact smooth data also follow
from \cite[Theorem A, equation (4)]{BahuaudHelliwell2015Uniqueness}
with \(k=1,\ c=1,\ a=-(1-\alpha)/6,\ b=(2+\alpha)/6\).
The strict symbol inequality is \(a>-1/6\). This application supplies
the reverse DeTurck step: the vector field in Bour's Proposition 2.1
depends on the fixed background and its pullback changes that background.

On a slab with bounded normalized curvature, the positive Yamabe bound
gives a Sobolev inequality with bounded lower-order term. A cutoff test
on balls gives a positive lower volume ratio at bounded scales, hence
injectivity at bounded curvature. Bour's Lemma 6.2 bounds the time
derivative of a curvature derivative energy \(\mathcal R_m\) plus
\(c_\alpha\mathcal R_{m+2}\) by
\(C_{m,\alpha}\|\Rm\|_\infty^{m+2}\mathcal R_0\).
In the subsequent time-weighted sum \(m=2j\), so the power on the right
is \(2j+2\). For \(K=\sup|\Rm|\) the resulting Gronwall coefficient
is bounded by \(C_{\alpha,k}K^2(1+K^2t)^k\), sufficient on
every bounded normalized slab. The resulting even-derivative
\(L^2\) bounds control the odd derivatives by integration by parts.
With the scalar Sobolev inequality, Kato's inequality gives \(L^4\)
bounds on every curvature derivative. Applying the same inequality
to the square of each tensor norm gives \(L^8\) bounds, and the
multiplicative Sobolev inequality \cite[Proposition 9.16]{Bour2010FourthOrderCurvatureFlows}
with \(q=m=8,\ p=\infty\) gives the pointwise bounds. This suffices for the positive-time
derivative estimates used here. These estimates are for fixed \(\alpha\).
The flow preserves volume since the functional is invariant under
constant metric scaling in dimension four. With bounded curvature,
the fixed-slab estimates give finite-time continuation, and in infinite
time there are times tending to infinity with gradient norm tending
to zero. Unit volume and ball noncollapse bound diameter. Compact
smooth convergence therefore gives a critical metric on the whole
original smooth type, up to diffeomorphism.

If curvature is unbounded, use first hits of increasing curvature
levels \(Q_i\to\infty\) and base at a maximum point at time \(t_i\).
On every fixed backward interval in \(s=Q_i^2(t-t_i)\) the metrics
\(Q_i g(t_i+s/Q_i^2)\) have bounded curvature, full curvature energy,
and the same Yamabe lower bound. The positive-time estimates starting
on a longer backward interval control derivatives through \(s=0\).
Pointed smooth compactness can use the markings of the final slice
and the evolution equation to compare on the same compact domains in
space and time: bounds on the flow tensor and its covariant derivatives
give uniform metric equivalence and bounds on the change of connection
by integrating \(\partial_s\Gamma=h^{-1}*\nabla\partial_s h\) backwards.
Differentiating in the terminal reference connection gives interior
bounds in every fixed spatial coordinate order. The equation gives
the time derivatives on a shorter closed slab through \(s=0\). The rescaled dissipation on a fixed backward interval
equals the dissipation on its original time interval. It tends to
zero because the total is finite:
\[
 \int_0^t\|\operatorname{grad}\F^\alpha\|_{L^2}^2\,ds
   =\tfrac12\big(\F^\alpha(g(0))-\F^\alpha(g(t))\big).
\]
The rescaled limit is stationary, including at the final slice by
smoothness. Its curvature norm equals one at the base point. It is
noncompact: a compact limit would embed onto all of \(M\), whereas
the rescaled volumes diverge. Its complete geometry has infinite
volume by the uniform lower bounds on unit balls. The trace equation
is \(\alpha\Delta R=0\). The inherited \(S\)-bound and a cutoff
test against \(R\) show that this harmonic function is constant,
hence zero. The remaining Euler equation is Bach-flatness.
Compactly supported Yamabe tests transport to the source, giving
\(Y_c(X,h)\ge y\). In particular \(y\|u\|_4^2\le6\|du\|_2^2\)
on this scalar-flat limit.

In either alternative, smooth comparison on compact sets preserves
the integrals there. Exhaustion bounds each of \(E_+,E_-,D,S\) and
the full curvature integral on the pointed limit separately. The
lower bound \(I\le e\) can be used again on a compactified limit only
after identifying its smooth type and positive class.

\section{Recovering the compact smooth type}
\begin{lemma}\label{lem:mark}
Suppose closed metrics on the oriented \(M\) have \(Y\ge y>0\),
\(E_\pm\le2T-\eta\), \(D\le d_*<32\pi^2\), and bounded full
curvature integral. Suppose constant rescalings converge in the smooth
pointed sense to a complete noncompact connected nonflat scalar-flat
Bach-flat \((X,h)\). Then \(X\) has one asymptotically Euclidean end.
Its smooth conformal compactification is orientation-preservingly
diffeomorphic to \(M\), has \(Y\ge y\), and has common chiral energy
satisfying
\[
   T<e_X\le\liminf e_i,\qquad 4e_X=\int_X|\Ric_h|^2\dV_h\le d_* .
\]
Here the liminf is along the converging subsequence.
\end{lemma}
\begin{proof}
Choose the limit orientation by the comparison at the base point,
passing to a subsequence for consistency. Use comparison embeddings
on a connected neighborhood of each \emph{whole} compact truncation
and both buffers on every end in that truncation. These embeddings
come from pointed smooth convergence on an exhaustion. Compact tests
and compact exhaustion give the homogeneous Sobolev inequality and
finite full curvature. Tian--Viaclovsky's decay theorem gives
\(|\nabla^m\Rm|=o(r^{-2-m})\) \cite[Theorem 1.1, first part]{TianViaclovsky2005ALE}.
The negative Ricci tensor is in \(L^2\), so the homogeneous Sobolev
inequality and this decay give finitely many order-zero ALE ends
\cite[Theorem 1.1]{TianViaclovsky2008Volume}. In particular the
end topology is finite and \(b_1(X)<\infty\), so the differentiated
order-zero charts in the first theorem are available.
Ache--Viaclovsky improve the scalar-flat Bach-flat ends to order two
with differentiated decay \cite[Theorem 1.2]{AcheViaclovsky2011Obstruction}.

On a large annulus of each end conformally multiply by \(r^{-2}\);
in logarithmic radius the metric approaches the round product on
\(S^3/\Gamma\). Fix buffer length, then radius, then comparison index.
Extend the positive change of gauge over the source using a cutoff
of its logarithm. A fixed smooth cutoff interpolation to the exact
product in a central subcollar has added \(E_\pm\) at most \(C\eta_2^2\)
per splice when the product \(C^2\) error is \(\eta_2\).
Use disjoint splices on the \emph{uncut} source. Their total can be
made arbitrarily small. In particular the resulting closed metric
has \(E_+<9\pi^2\) and is exactly product at all selected slices.
The closed-source quotient obstruction \cite[A quotient obstruction on the closed source]{Baker2026MarkedNecks}
applies to each. It uses Hillman's finite-group restriction, the spin
obstruction for the binary icosahedral link, and the product-boundary
signature contribution \(9\pi^2\) on one side of a quaternionic link.
Thus all end groups are trivial.

Let \(k\) be the number of ends. The embedded connected core with its
\(k\) disjoint boundary spheres has one exterior component per sphere:
each sphere separates in \(M\), and the incidence graph is a tree.
Mayer--Vietoris and van Kampen show that every capped piece is simply
connected with \(H_2=0\), hence a homotopy sphere. Removing \(k\) balls
from the central capped piece gives \(\chi(X)=2-k\). In Gauss--Bonnet
each outgoing Euclidean end contributes \(8\pi^2\), by convergence of
the rescaled boundary metric and second fundamental form to those of
the round unit sphere \cite[\S7]{TianViaclovsky2005ALE}. Hence
\[
 D_X=2\big(E_+(h)+E_-(h)\big)+32\pi^2(k-1).
\]
The bound \(D_X\le d_*\) forces \(k=1\).

On that end choose a positive harmonic function
\(a=r^{-2}(1+O(r^{-\beta}))\), with the expansion differentiated
through order five for \(0<\beta<1\). For example solve exterior
Dirichlet problems on \(r_0\le r\le L\) with boundary values
\(r_0^{-2}\) and \(L^{-2}\), taking \(r_0\) sufficiently large.
Small and large positive multiples
of \(r^{-2}(1+r^{-\beta})\) and \(r^{-2}(1-r^{-\beta})\) give lower
and upper barriers at sufficiently large radius. The resulting
two-sided comparison with \(r^{-2}\) and interior elliptic estimates
give \(\Delta_{\rm Euc}a=O(r^{-6})\) in the ALE chart. Cut off at
the inner boundary and expand the Euclidean Newton potential using
its finite \(\beta\)-moment. In the region \(|z|\le r/2\), the kernel
difference from \(r^{-2}\) is bounded by \(C r^{-2-\beta}|z|^\beta\).
On the tail, split at \(|z-x|=r/4\) and use the \(O(|z|^{-6})\)
source bound near \(x\) and the integrable tail mass elsewhere.
The decaying entire harmonic difference is zero.
The leading coefficient of \(r^{-2}\)
is positive by the lower barrier. After normalization, rescaled interior estimates give the differentiated remainder. Extend \(a\) positively
across the core. After inversion \(a^2h\) is scalar-flat near the
puncture, Bach-flat there and \(O(\rho^\beta)\) close to Euclidean
with \(m\)th derivative error \(O(\rho^{\beta-m})\).
The removable point theorem \cite[Definition 1.6, Theorem 1.8 and
Remark 7.2]{AcheViaclovsky2011Obstruction} gives a smooth chart
for the compactification \(N\). Bach-flatness holds on the punctured
manifold by conformal covariance and extends across the point by
smoothness. Its topology is the one-point
compactification of the one-ended interior, hence that of a homotopy
sphere. A point has zero \(H^1\) capacity in dimension four.
Transport compact tests and cut off near the point to obtain
\(Y(N,[a^2h])=Y_c(X,h)\ge y\). Conformal invariance gives both
energies, equal by the signature of \(N\), whence \(D_X=4e_X\).
If the compactification were conformally round then \(e_X=D_X=0\)
would give \(\Rm_h=0\). The closed critical gap on a homotopy sphere
therefore gives \(e_X>T\)
\cite{ChangGurskyYang2003,ChangChen2021RicciSphere,Baker2026MarkedNecks}.

For the smooth identification use a small \emph{normal} ball of radius
\(s\) in the removable chart of \(N\). Its complement and both
logarithmic buffers are contained in a single larger compact truncation
of \(X\). Under the same embedding of this truncation, multiplying the
pulled-back, constantly rescaled source by \(a^2\rho^{-2}\) gives
product error \(O_\ell(s^2)+o_i(1;s,\ell)\) in \(C^2\) around
the normal sphere of radius \(\rho=s\). If \(e_i\) has no limit,
pass to a subsequence with convergent \(e_i\), or use the common
upper bound. The raw exterior energy has upper limit at most
\(2T-\eta-e_X+E_+(N;B_s)<T\) for small \(s\).
Choose \(\ell\) first to absorb the \(O(\ell^{-3/2})\) Yamabe
localization loss in the buffered cap theorem
\cite[One buffered cut and finite ordinary punctures]{Baker2026MarkedNecks}.
Then choose \(s\) small and the comparison index large to retain
positive exterior Yamabe sign and capped energy strictly less than
\(T\), including the splice cost. The exterior capped by the actual
parametrized ball is diffeomorphic to \(S^4\) by positive recognition.

Double the removed normal ball in \(N\). The capped source exterior
and this double are recognized oriented standard spheres. Send the
actual cap and a collar to the opposite ball in the double by the
given collar coordinates, with the prescribed inverse comparison on
the adjoining strip. In collar polar coordinates the ball map can
be made linear near its center. The disk-neighborhood extension
theorem \cite[Theorem C]{Palais1960Extending} realizes this local
identification by an oriented diffeomorphism of the two spheres,
after shrinking the collar if needed. Its restriction identifies the
source exterior with the actual deleted normal ball and agrees with
the inverse core comparison near the seam. Gluing these maps gives
\(M\to N\), and the two inverse restrictions give \(N\to M\).
Lastly \(E_\pm(h)\le\liminf E_\pm(g_i)\) by compact exhaustion.
More precisely, if \(e_i\to e_{\rm src}\), the limiting raw exterior
mass at fixed \(s\) is \(e_{\rm src}-e_X+E_\pm(N;B_s)\) for each
sign, by subtraction on the original closed metric. Only after this
limit do we let \(s\downarrow0\).
\end{proof}

\paragraph{Restriction on fixed-flow bubbling.}
If one of the late original flows has a noncompact model, the marked
identification gives \(I\le e_X\). Keeping the exact \(D_0\) in
Proposition~\ref{prop:budget} and using \(D_X=4e_X\) gives
\[
  4(e_X-I)+\frac{\Yrd^2-Y_0^2}{12}
     \le\frac{4(e_0-I)}{\alpha}.
\]
Thus a subsequence of such original flows can bubble only if its
\emph{initial} Yamabe constants satisfy
\(\Yrd^2-Y_{j,0}^2\le48\alpha_j\). The iterated upper limit of raw exterior mass in
either chirality is at most \(e_{j,0}/(1-\alpha_j)-I\) at fixed \(j\),
and hence tends to zero along that subsequence.

\section{The limit of critical metrics}
At each late \(j\), choose the closed critical metric of Bour's
alternative if this is the output. Keep its physical volume-one
representative. If the output is pointed, apply Lemma~\ref{lem:mark}
to the source slices, pull back the exact compact class and choose
its volume-one Yamabe representative. Write \(p_j=\alpha_j\) in
the first case and \(p_j=0\) in the second. These closed metrics
\(q_j\) all lie on \(M\). They have \(I\le e(q_j)\le e_{j,0}/(1-\alpha_j)\),
so \(e(q_j)\to I\). In the second case we recompute \(D(q_j)\le4e(q_j)\)
in Yamabe gauge. In the first case Proposition~\ref{prop:budget}
passes through whole-manifold smooth convergence. In both cases
\(D(q_j)\le d_*<32\pi^2\) eventually and the full curvature integral
is uniformly bounded.

All these metrics are constant-scalar. For a regularized critical
point the trace gives \(\Delta R=0\) since \(p_j>0\).
The unit-volume constant test and scalar square bound give
\(y\le R_j\le C\); the same bounds hold in the exact Yamabe case.
Write \(E=\Ric_0\), use Bour's nonnegative rough Laplacian and
\(\mathring W(E)\) for his curvature action. On these representatives
the trace-free Euler equation is
\[
 0=\tfrac12\nabla^*\nabla E-\mathring W(E)
       +E\circ E-\tfrac14|E|^2q_j+\tfrac{2-p_j}{12}R_j E
 \quad\text{\cite[Proposition 4.4]{Bour2010FourthOrderCurvatureFlows}.}
\]
In particular \(\operatorname{grad}\F_W=p_j R_j E/12\) in this
convention, and its vanishing is Bach-flatness. The principal
coefficient is independent of \(p_j\).

For \(g=\lambda q_j\), \(\lambda\ge1\), the scalar curvature is
\(R_j/\lambda\) and
\[
 y\|u\|_{4,g}^2\le6\|du\|_{2,g}^2+(R_j/\lambda)\|u\|_{2,g}^2 .
\]
Testing a cutoff of \(B(x,r/2)\) in \(B(x,r)\) gives
\(y V(r/2)^{1/2}\le(24/r^2+C/\lambda)V(r)\).
Iterating on halved radii gives \(V(r)\ge c r^4\) for \(0<r\le1\),
with \(c>0\) uniform. Volume one then bounds the unscaled diameter.
Choose \(\lambda=1\) on a subsequence where \(\sup|\Rm|\) is bounded;
otherwise take \(\lambda_j=\max|\Rm(q_j)|\to\infty\), based at a
maximum. In either normalized gauge curvature is bounded and the
injectivity radius has a uniform positive lower bound by
\cite[Theorem 4.7]{CheegerGromovTaylor1982}.

Bounded sectional curvature and injectivity give harmonic charts of
uniform \(C^{1,\gamma}\) size, \(0<\gamma<1\)
\cite[Theorem 4.3]{Peters1987Convergence}.
For clarity, the interior elliptic bootstrap in these charts uses
no bound on a derivative of the scalar. The harmonic-coordinate
Ricci identity gives uniform \(W^{2,p}\) control of the metric \cite[Lemma 4.1 and (4.3)]{DeTurckKazdan1981Regularity} for
fixed \(p>4\) on smaller balls. In the rough Laplacian of \(E\),
\((\partial\Gamma)E\) is therefore \(L^p\), and \(\Gamma\partial E\)
is a bounded first-order coefficient. The displayed Euler equation
has uniformly bounded algebraic right side and gives a \(W^{2,p}\)
bound for \(E\) by interior estimates with lower-order terms.
Thus \(E\) is \(C^{1,\gamma'}\) for some \(\gamma'>0\).
Alternating the metric equation and the \(E\) equation now gives
uniform interior estimates of every order. Both equations have the
same form with the rescaled tensors, with \(R_j/\lambda\) in the
zeroth-order term. Smooth pointed compactness gives one comparison
on every connected compact domain.

At bounded unscaled curvature the diameter bound gives a closed
limit and whole oriented comparison diffeomorphisms. It has \(Y\ge y\)
by testing smooth convergence. Since \(p_j\to0\), its limiting Euler
equation is Bach-flat and smooth convergence gives \(e=I\).
At unbounded curvature the complete smooth limit has curvature norm
one at its base point. The volumes of the rescaled metrics tend to
infinity, precluding a compact limit with a whole comparison
diffeomorphism. The scalar vanishes since \(R_j/\lambda_j\to0\);
the Euler equation gives Bach-flatness. Lemma~\ref{lem:mark} applies
to this second pointed sequence, giving a positive smooth compactification
on \(M\) with
\[
 I\le e_X\le\liminf_j e(q_j)=I .
\]
In this second pointed case the iterated raw exterior mass is
\(I-e_X=0\) for both signs. This finishes both alternatives without
a restart. On a nonstandard
type, positive recognition below \(T\) excludes \(I=0\). The attained
value thus belongs to the set of positive subcritical exact values.
Its infimum \(a_M\) satisfies \(I\le a_M\le I\). Finally positive Yamabe
sign is open under smooth variations, so the usual second derivative
test at the minimizing metric proves the stability assertion of
Theorem~\ref{thm:main}.

\paragraph{Scope and provenance.}
Drafted for Edward B. Baker III with OpenAI Codex assistance,
26 November 2026, from checkout \texttt{4da5865}. The starting proposal
was F2 of the author's Claude review; the fixed-parameter source is
Bour v1. The selected research configuration was GPT-6.2 Astra/ultra
with bounded GPT-6.2 Sol/ultra calculations and readings. Serving
telemetry is unavailable. The argument uses the cited literature and
the project's marked-neck working results. The present theorem leaves
the construction of positive subcritical entry on an arbitrary type
and the exclusion of an attained nonround minimum on a nonstandard type
open.

\bibliographystyle{unsrt}
\bibliography{references}
\end{document}
